\documentclass[9pt,a4paper]{extarticle}

\usepackage[T1]{fontenc}
\usepackage[utf8]{inputenc}
\usepackage[brazil]{babel}
\usepackage{geometry}
\geometry{margin=2cm}
\usepackage{parskip}
\usepackage{setspace}
\usepackage[protrusion=true,expansion=false]{microtype}
\usepackage{enumitem}
\usepackage{amsmath, amssymb, amsthm, bm}
\usepackage{algorithm}
\usepackage{algpseudocode}
\usepackage{xcolor}
\usepackage{booktabs}
\usepackage{array}
\usepackage{tabularx}
\usepackage{float}
\usepackage{hyperref}
\hypersetup{colorlinks=true, linkcolor=blue, urlcolor=blue}

\definecolor{important}{RGB}{200, 50, 50}
\definecolor{notice}{RGB}{0, 102, 204}
\newcommand{\impt}[1]{\textcolor{important}{\textbf{#1}}}
\newcommand{\ntc}[1]{\textcolor{notice}{#1}}

\begin{document}

\begin{center}
  \Large \textbf{Prova 3 \\ Estruturas de Dados Elementares e Árvores} \\[4pt]
  \normalsize Disciplina: PCC104 -- Projeto e Análise de Algoritmos \quad|\quad DECOM -- UFOP \\[2pt]
  \normalsize \today
\end{center}

\medskip
\noindent\textbf{Instruções.} Esta prova contém 7 questões discursivas. Você deverá entregar o código das implementações da lista de exercícios impresso juntamente com a prova. Algumas das questões farão referência ao seu próprio código. A não entrega do código implica em nota 0 nestas questões.

\bigskip\hrule\bigskip

\begin{enumerate}[label=\textbf{\arabic*.}, leftmargin=*, itemsep=2em]

% ------------------------------------------------------------------
% Q1 – Arrays
% ------------------------------------------------------------------
\item \textbf{(Arrays.)}
  Escreva o pseudocódigo de um algoritmo que recebe um array \(A[1\,..\,n]\) já
  \emph{ordenado em ordem crescente} e um valor \(x\), e devolve o índice de
  \(x\) em \(A\) (ou \textsc{nil} caso \(x\) não esteja presente). Apresente a análise de complexidade deste algoritmo.
  
% ------------------------------------------------------------------
% Q2 – Pilhas
% ------------------------------------------------------------------
\item \textbf{(Pilha com mínimo em \(O(1)\).)}
  Considere sua implementação de \textsc{MinStack}.
  Execute a sequência de operações abaixo sobre uma instância inicialmente vazia,
  mostrando após cada chamada o estado das duas pilhas e o valor devolvido
  (quando houver):
  \[
    \textsc{Push}(8),\;
    \textsc{Push}(3),\;
    \textsc{Push}(5),\;
    \textsc{Push}(1),\;
    \textsc{Min}(),\;
    \textsc{Pop}(),\;
    \textsc{Min}(),\;
    \textsc{Pop}(),\;
    \textsc{Min}()
  \]
  Em seguida, enuncie o invariante que sua estrutura auxiliar mantém e explique
  por que \textsc{Min}() é \(O(1)\) no pior caso.
% ------------------------------------------------------------------
% Q3 – Filas
% ------------------------------------------------------------------
\item \textbf{(Filas.)}
  Escreva o pseudocódigo de \textsc{Enqueue}$(Q, x)$ e \textsc{Dequeue}$(Q)$ para
  uma fila implementada em um array circular \(Q[0\,..\,m-1]\) com atributos
  \(Q.\mathit{head}\), \(Q.\mathit{tail}\) e \(Q.\mathit{size}\). Indique o custo de cada operação.

% ------------------------------------------------------------------
% Q4 – Listas duplamente encadeadas
% ------------------------------------------------------------------
\item \textbf{(Listas duplamente encadeadas.)}
  Seja \(L\) uma lista duplamente encadeada com sentinela \(L.\mathit{nil}\) no
  estilo CLRS, contendo \(n\) elementos.
  Justifique a complexidade de tempo de \textsc{List-Insert} (inserção no início),
  \textsc{List-Delete} (dado ponteiro \(x\) para o nó) e \textsc{List-Search}
  (busca por chave \(k\)).

% ------------------------------------------------------------------
% Q5 – Tabelas hash
% ------------------------------------------------------------------
% ------------------------------------------------------------------
% Q5 – Tabelas hash
% ------------------------------------------------------------------
\item \textbf{(Tabelas hash.)} O pseudo-código abaixo apresenta o algoritmo de inserção em uma tabela hash com encadeamento separado, onde cada posição aponta para uma lista encadeada e a função hash é dada por \(h(k) = k \bmod m\).

  \begin{algorithmic}[1]
  \Procedure{Hash-Insert}{$T, k$}
    \State $j \gets k \bmod m$
    \State cria novo nó $z$ com $z.\mathit{key} = k$
    \State $z.\mathit{next} \gets T[j]$
    \State $T[j] \gets z$
  \EndProcedure
\end{algorithmic}

  Execute \textsc{Hash-Insert} para cada chave da sequência
  \[
    10,\; 22,\; 31,\; 4,\; 15,\; 28,\; 17
  \]
  em uma tabela com \(m = 7\), mostrando o estado completo de \(T\) após cada inserção.

% ------------------------------------------------------------------
% Q7 – Altura de ABB  (mantida como estava)
% ------------------------------------------------------------------
\item \textbf{(Altura de uma árvore binária de busca.)}

  \begin{enumerate}[label=(\alph*)]
    \item Escreva o pseudocódigo de \textsc{Tree-Height}$(x)$, que recebe a raiz
          \(x\) de uma ABB e devolve sua altura (árvore vazia $= -1$, nó folha $= 0$).
          Apresente a versão recursiva, identificando caso base e passo recursivo.
    \item Formule a recorrência \(T(n)\) no pior caso, resolva-a e expresse o
          resultado em notação \(\Theta\).
    \item Uma ABB é construída inserindo as chaves \(10, 5, 15, 3, 7, 12, 20\)
          nessa ordem. Desenhe a árvore e aplique \textsc{Tree-Height} à raiz,
          mostrando os valores devolvidos em cada chamada recursiva.
  \end{enumerate}

% ------------------------------------------------------------------
% Q8 – Árvores rubro-negras (simplificada)
% ------------------------------------------------------------------
\item \textbf{(Árvores rubro-negras.)}
  Quando devemos utilizar árvores rubro-negras? 
\end{enumerate}

\end{document}